% ============================================================
%  2. IDENTIFICACIÓN DEL PROBLEMA
%  Propósito: identificar y delimitar con precisión la situación
%  problemática que da origen al proyecto.
%  CUIDADO (regla explícita de la guía): no confundir el problema
%  con la solución. "Implementar Random Forest" NO es un problema:
%  es una técnica que podría evaluarse después.
% ============================================================
\chapter{Identificación del Problema}
\label{cap:problema}

\section{Descripción del problema}
\label{sec:descripcion-problema}

Trufi Association coordina el relevamiento colaborativo de las rutas de transporte público
del eje metropolitano de Cochabamba. Su labor consiste en recorrer las líneas de micros y
trufis, registrar sus trayectos y paradas, verificarlos con los operadores y convertirlos
al formato GTFS que alimenta el planificador de ``Trufi App''. Esta tarea exige
desplazamientos físicos, tiempo de procesamiento técnico y coordinación entre personas que
participan de forma voluntaria. Por consiguiente, la capacidad humana disponible limita la
cantidad de zonas que pueden mapearse o actualizarse en un período determinado, y cada
jornada de trabajo debe dirigirse a alguna zona en particular. La unidad de análisis
adoptada es la zona hexagonal de la grilla geoespacial H3 (resolución 8), con un área
aproximada de 0,74 km². En el resto del documento se la denomina indistintamente ``zona
hexagonal'' o ``zona''.

En la situación actual, la elección de las zonas a relevar se apoya principalmente en el
conocimiento acumulado del equipo y en observaciones informales del territorio. La
organización dispone de un registro detallado de las consultas realizadas en la aplicación
a lo largo de casi dos años, pero dicho registro no se aprovecha de manera sistemática para
orientar la asignación del trabajo voluntario.

La dificultad no radica en la ausencia de datos, sino en su interpretación. El número de
consultas que genera una zona depende simultáneamente de varios factores: la población
residente, su localización respecto de los centros de actividad y la cobertura de la red
mapeada, es decir, la utilidad que la aplicación ofrece a sus habitantes. Mientras estos factores no se separen, un conteo bajo resulta
ambiguo, puesto que puede reflejar tanto una demanda naturalmente reducida como una
necesidad no atendida.

Las consecuencias de esta ambigüedad recaen, en última instancia, sobre los usuarios del
transporte. Una zona con rutas no mapeadas no aparece en los itinerarios que ofrece la
aplicación, de modo que sus habitantes no pueden planificar viajes que quizá sí son
posibles, o deben recurrir a alternativas más costosas o más lentas. La brecha de
información se traduce así en viajes no cubiertos o mal planificados, y en una percepción
de que la aplicación no resulta útil precisamente donde el transporte informal es la
principal opción de movilidad. En el plano operativo, la ambigüedad también tiene un costo:
ante la imposibilidad de distinguir una demanda baja de una necesidad no atendida, el
esfuerzo de mapeo se distribuye sin un criterio cuantitativo que permita concentrarlo
donde el impacto potencial es mayor.

La evidencia disponible sugiere que esta ambigüedad es relevante. El antecedente directo de
este trabajo, desarrollado con registros de la misma aplicación, encontró que la mayor
parte de las zonas analizadas debía excluirse de la predicción por no contar con datos
suficientes, y que la demanda se concentraba en unas pocas zonas del territorio (Angulo
Andrade, 2024). En consecuencia, las zonas con menor actividad, que podrían ser
precisamente aquellas donde la aplicación resulta menos útil, son también las que los
enfoques orientados a pronosticar el volumen de consultas dejan fuera del análisis.

El problema se delimita a la variación de la demanda de información entre zonas del eje
metropolitano, y no a su evolución en el tiempo. Desde la perspectiva de la Ciencia de
Datos, lo que se busca apoyar es una estimación: el número de consultas que cabría esperar
en cada zona según su población residente y su localización —la distancia al centro y el
comportamiento de las zonas vecinas—, junto con una medida de la confiabilidad de dicha
estimación en zonas no observadas. La cobertura de la red mapeada no interviene en el
cálculo de lo esperado: se reserva para interpretar la diferencia entre lo observado y lo
esperado, porque si se usara para estimar las consultas esperadas, una zona sin rutas
mapeadas tendría, por construcción, una expectativa baja, y la brecha que se pretende
detectar quedaría absorbida por el propio modelo. La estimación corresponde al
total de consultas acumulado en cada zona a lo largo del período analizado, de modo que
el conteo esperado de una zona hexagonal pueda compararse con el total de consultas
efectivamente registradas en ella entre septiembre de 2022 y junio de 2024. Con esta
referencia, la diferencia entre lo observado y lo esperado podría emplearse como criterio
cuantitativo para priorizar la revisión del mapeo.

\section{Formulación del problema}
\label{sec:formulacion-problema}

\begin{quote}
  \textbf{¿Cómo estimar las consultas esperadas de ``Trufi App'' por zona hexagonal H3 en el eje
   metropolitano de Cochabamba, a partir de su población y su entorno espacial, 
   para orientar la priorización del mapeo de rutas?}
\end{quote}
